\section{The global comparison metric}\label{sec:reconstruction}

Return to the family of Section~\ref{sec:geometry}. All parameters
are constant during a trajectory. Let $H(s)$ solve
\eqref{eq:finite-flow} in a fixed product ball about $o$. Put
\begin{equation}\label{eq:small-parameters}
 \alpha_* =\max(|\alpha_+|,|\alpha_-|),\qquad
 \rho=\epsilon_*,\qquad \ell=1+|\log\rho|,\qquad
 \sigma=\alpha_*^2+\epsilon_*.
\end{equation}
We always take $\alpha_*\to0$ and $L_*\to\infty$, with no
restriction on their ratios or the individual $L_e$. Constants in
this section depend on the fixed component data and product ball.
They can be chosen uniformly over unitary transports. The physical
speed is $O(\sigma)$; its product with $L_*$ need not be bounded.

For later estimates define
\begin{align}
 d_e&=d(H_i,U_{ij}H_jU_{ij}^*),&
 F_{\mathrm{ed}}&=\sum_e\epsilon_e d_e,&
 E_{\mathrm{ed}}&=\sum_e\epsilon_e d_e^2,\label{eq:edge-quantities}\\
 m_i&=\|[T_i,H_i]\|,&
 F_{\mathrm{co}}&=\sum_{i=+,-}\alpha_i^2\kappa_i
                            \sqrt{2f_{T_i}(H_i)},&
 V&=\|\dot H\|_{g_A}.\label{eq:component-quantities}
\end{align}
Write $F=F_{\mathrm{ed}}+F_{\mathrm{co}}$. Fixed factors $2\pi$
are absorbed in constants. Smoothness on the bounded matrix range and
\eqref{eq:cosh-potential} give
\begin{equation}\label{eq:force-energy-relations}
 \begin{gathered}
 V\le CF,\qquad F_{\mathrm{ed}}\le C\epsilon_*,
 \qquad |\dot B_{ij}|\le CV,\\
 F_{\mathrm{ed}}^2\le C\epsilon_*E_{\mathrm{ed}},\qquad
 \sum_i\alpha_i^2m_i^2+E_{\mathrm{ed}}\le C\Phi(H),
 \qquad \sum_i\alpha_i^2m_i\le CF_{\mathrm{co}}.
 \end{gathered}
\end{equation}
For example $[T,e^S]=\int_0^1e^{(1-t)S}[T,S]e^{tS}\,dt$
bounds $m_i$ by $C\|[T_i,\log H_i]\|$. The finite number of
edges gives the first inequality on the second line by
Cauchy--Schwarz.

\subsection{Component Green functions and node profiles}
Let $g_{ie}$ be the mean-zero Green function on the fixed component
$C_i$ with
\begin{equation}\label{eq:green}
 \Delta_i g_{ie}=-2\pi\delta_{p_{ie}}+2\pi/A_i,
 \qquad g_{ie}(u)=\log|u|+\gamma_{ie}(u).
\end{equation}
In the harmonic unitary gauge put
\begin{equation}\label{eq:component-correction}
 \xi_i=-\sum_{e=ij}\epsilon_e B_{ij}g_{ie},
 \qquad K_i=H_i e^{\xi_i},\qquad
 \Delta_i\xi_i=-\frac{2\pi}{A_i}\sum_{e=ij}\epsilon_e B_{ij}
\end{equation}
off the punctures. The matrices $\xi_i$ are $H_i$-self-adjoint
and trace-free. On the component cut off at radius $\rho/2$, the
fixed Green functions give
\begin{equation}\label{eq:green-bounds}
 \begin{gathered}
 \|\xi_i\|_\infty\le C\ell F_{\mathrm{ed}},\qquad
 \int(|\xi_i|+|\nabla\xi_i|)\le CF_{\mathrm{ed}},\qquad
 \int|\xi_i|^2\le CF_{\mathrm{ed}}^2,\\
 \int|\nabla\xi_i|^2\le C\ell F_{\mathrm{ed}}^2,
 \qquad \int|\xi_i||\Delta_i\xi_i|\le CF_{\mathrm{ed}}^2,
 \qquad \int|\partial_s\xi_i|\le C\epsilon_*V.
 \end{gathered}
\end{equation}
The logarithmic loss comes only from the integral of $|\nabla\log|u||^2$.

At incidence $ie$ expand
\begin{equation}\label{eq:regular-split}
 \xi_i(u)=-\epsilon_e B_{ij}\log|u|+c_{ie}+q_{ie}(u),
 \qquad q_{ie}(0)=0,\qquad P_{ie}=H_i e^{c_{ie}}.
\end{equation}
The constant $c_{ie}$ is the regular part of the complete Green sum,
including all other punctures of the component. It satisfies
$|c_{ie}|\le CF_{\mathrm{ed}}$, and on
$\rho\le|u|\le2\rho$,
\begin{equation}\label{eq:regular-bounds}
 |\nabla^kq_{ie}|\le CF_{\mathrm{ed}}\rho^{1-k}
 \quad(0\le k\le2).
\end{equation}
The corresponding time derivatives have $F_{\mathrm{ed}}$ replaced
by $\epsilon_*V$.

In the node frame at its $i$ end define
\begin{equation}\label{eq:node-profile}
 N_e(s,u)=\Geo(P_{ie},U_{ij}P_{je}U_{ij}^*;\theta_e),
 \qquad \theta_e=-\epsilon_e\log|u|.
\end{equation}
Here $\Geo(P,Q;\theta)=P\exp(\theta\log(P^{-1}Q))$.
The expression from the opposite end gives the same metric: the
parameter becomes $1-\theta_e$ because $uv=\tau_e$ and
$\epsilon_e L_e=1$. This identity requires the individual
$\epsilon_e$, rather than the largest $\epsilon_*$. For each
fixed time, $N_e^{-1}\partial N_e$ is a constant matrix times
$\partial\log|u|$, so its Chern curvature vanishes on the annulus.
Its endpoint velocities are at most $CV$.

Let $\chi(|u|/\rho)$ be a smooth radial cutoff, zero near the inner
circle of the collar $\rho<|u|<2\rho$ and one near the outer circle.
The component metric in the same flat node frame is $G_{ie}^*K_iG_{ie}$.
On the collar set
\begin{equation}\label{eq:collar-metric}
 \widehat K=H_i\exp\bigl((1-\chi)\log(H_i^{-1}N_e)
             +\chi\log(H_i^{-1}G_{ie}^*K_iG_{ie})\bigr).
\end{equation}
Use $K_i$ on the component exterior and $N_e$ on the middle node.
All logarithms in \eqref{eq:collar-metric} are $H_i$-self-adjoint
and trace-free. These formulas therefore give a smooth positive metric
on the actual bundle, with the same determinant as $h_0$. They are
equivariant under the constant unitary node transitions. Since
$\epsilon_*\ell\to0$, the component corrections are small and
the comparison metrics have fixed two-sided bounds for all times.
On every fixed compact component region,
\begin{equation}\label{eq:compact-reconstruction}
 d(\widehat K(H),H_i)\le C\epsilon_*.
\end{equation}

\subsection{The full component curvature calculation}
For spatially constant $H_i$ in the harmonic gauge, put
$B_i=H_i^{-1}T_iH_i$. The Chern connection satisfies
\[
 H_i^{-1}\partial_{D_i}H_i
       =-i\alpha_i(B_i-T_i)\vartheta_i^{1,0},\qquad
 F_{H_i}=\alpha_i^2[T_i,B_i]
                  \vartheta_i^{1,0}\wedge\vartheta_i^{0,1}.
\]
With our metric convention,
$2i\Lambda(\vartheta_i^{1,0}\wedge\vartheta_i^{0,1})
=|\vartheta_i|^2=\kappa_i/A_i$. Thus
\begin{equation}\label{eq:constant-curvature}
 \mathcal C_i(H_i):=2i\Lambda_iF_{H_i}
      =\frac{\alpha_i^2\kappa_i}{A_i}
                  [T_i,H_i^{-1}T_iH_i].
\end{equation}
Equations \eqref{eq:finite-flow} and \eqref{eq:component-correction}
give the exact leading cancellation
\begin{equation}\label{eq:component-cancellation}
 H_i^{-1}\dot H_i+\mathcal C_i(H_i)+\Delta_i\xi_i=0.
\end{equation}

We record how the complete Chern expression is controlled, since its
ordered terms are needed on an unbounded time interval. The identity
\[
 K_i^{-1}\partial_{D_i}K_i
  =e^{-\xi_i}(H_i^{-1}\partial_{D_i}H_i)e^{\xi_i}
                         +e^{-\xi_i}\partial_{D_i}e^{\xi_i}
\]
and the exact formula
\[
 e^{-\xi}\partial e^\xi
       =\int_0^1 e^{-t\xi}(\partial\xi)e^{t\xi}\,dt
\]
imply, after covariant differentiation,
\begin{equation}\label{eq:chern-remainder}
 \begin{split}
 |2i\Lambda_iF_{K_i}-\mathcal C_i(H_i)-\Delta_i\xi_i|
 \le C\bigl(&\alpha_*|\nabla\xi_i|+\alpha_*^2|\xi_i|
              +|\nabla\xi_i|^2+|\xi_i||\Delta_i\xi_i|\bigr).
 \end{split}
\end{equation}
Indeed the first two errors have respectively one connection
coefficient and a derivative of $\xi_i$, or two connection
coefficients and $\xi_i$. The coefficients of $\vartheta_i$ are
constant in flat coordinates. The remaining errors come from the
exponential integral and contain two derivatives of $\xi_i$ or its
product with its Laplacian. Further factors of $\xi_i$ are absorbed
by its small supremum norm. On rational components $\alpha_i=0$.
The time identity bounds
$|K_i^{-1}\partial_sK_i-H_i^{-1}\dot H_i|$ by
$C(|\partial_s\xi_i|+V|\xi_i|)$.

Combining these statements with \eqref{eq:green-bounds} gives the
component residual bound
\begin{equation}\label{eq:component-residual}
 \begin{split}
 \int|\cR(K_i)|\le C\bigl(&\epsilon_*V+VF_{\mathrm{ed}}
                 +\ell F_{\mathrm{ed}}^2
                 +\alpha_*F_{\mathrm{ed}}
                 +\alpha_*^2F_{\mathrm{ed}}\bigr)\\
       &\le C(\alpha_*+\epsilon_*\ell)(F+V).
 \end{split}
\end{equation}
The branch identification preserves both the complex structure and
the reference connection exactly. Transfer to the fibre changes only
the metric contraction and area. Lemma~\ref{lem:geometric-estimates}
adds at most
\begin{equation}\label{eq:transfer-residual}
 Cb^2\rho^{-N}(F+V)
\end{equation}
for one fixed finite $N$. To retain this force factor, note that the
constant-metric curvature in \eqref{eq:constant-curvature} is
bounded by $C\alpha_i^2m_i$; every other spatial term contains
$\xi_i$ or its derivatives. The latter lose only fixed powers of
$\rho^{-1}$. The middle-node time residual has integral at most
$C\rho^2V$, since its spatial curvature vanishes and its area is
$O(\rho^2)$. These estimates already have the desired integrable
form on a small finite basin. The incidence collars require the
additional argument in the next section.

\subsection{Energy and dependence on the finite tuple}
The comparison metric has an integrated energy bound even without
the small-ball force estimate.

\begin{lemma}\label{lem:reconstruction-energy}
On every fixed finite ball, for sufficiently small joint parameters,
\begin{equation}\label{eq:reconstruction-energy}
 \int_{C_\tau}|D\widehat K(H)|^2\le C\Phi(H),\qquad
 \int_0^\infty\!\int_{C_\tau}|D\widehat K(H(s))|^2
                     \le C d_A(H(0),o)^2.
\end{equation}
Also the map $H\mapsto\widehat K(H)$ is uniformly Lipschitz from
that ball to the integrated distance on $C_\tau$:
\begin{equation}\label{eq:reconstruction-lipschitz}
 \int_{C_\tau}d(\widehat K(H),\widehat K(H'))^2
                                      \le C d_A(H,H')^2.
\end{equation}
Constants are uniform in the unitary transports and all scale ratios.
\end{lemma}

\begin{proof}
The component derivative is covariant with respect to
$D_i=d-i\alpha_i[T_i,\cdot]\vartheta_i$. From
\eqref{eq:green-bounds} and bounded differentials of the exponential,
\begin{equation}\label{eq:component-energy}
 \int|D_iK_i|^2\le C\bigl(\alpha_i^2m_i^2
                       +(\ell+\alpha_i^2)F_{\mathrm{ed}}^2\bigr).
\end{equation}
Real one-form energy on a Riemann surface is conformally invariant,
so the exact holomorphic identification transfers this inequality
to the fibre without a metric error. Its sum is bounded by $C\Phi$
using \eqref{eq:force-energy-relations} and $\epsilon_*\ell\to0$.

The two corrected endpoints at edge $e$ have distance at most
$d_e+CF_{\mathrm{ed}}$. Constant geodesic speed in
\eqref{eq:node-profile} gives middle-node energy at most
\[
 C\epsilon_e^2(L_e+2\log\rho)(d_e+F_{\mathrm{ed}})^2
       \le C\epsilon_e(d_e+F_{\mathrm{ed}})^2.
\]
Its sum is again bounded by $C\Phi$.

On a collar, let $\delta_{ie}$ be the difference of the two
relative logarithms in \eqref{eq:collar-metric}. Differentiating
$G_{ie}^*H_iG_{ie}$ first retains the factor $m_i$ in its first
derivative. The action on $\xi_i$ contributes the other gauge
term below. The regular Green part and the mixed endpoint/geodesic
expansion give
\begin{equation}\label{eq:collar-discrepancy}
 \begin{split}
 |\delta_{ie}|+\rho|\nabla\delta_{ie}|&\le C a_{ie},\\
 a_{ie}&=|\alpha_i|\rho m_i+
       |\alpha_i|\rho\ell F_{\mathrm{ed}}+
       \rho F_{\mathrm{ed}}+\epsilon_e\ell F_{\mathrm{ed}}.
 \end{split}
\end{equation}
For the last term, at zero endpoint corrections the node logarithm
is exactly $\theta_e B_{ij}$; at $\theta_e=0$ it is $c_{ie}$.
The remainder from their sum vanishes on both axes and is at most
$C\theta_e F_{\mathrm{ed}}$, with the corresponding derivative
bounds. Since $\theta_e\le C\epsilon_e\ell$ on the collar,
this gives the displayed term. The derivative of the radial node
logarithm is at most
$C\rho^{-1}\epsilon_e(d_e+F_{\mathrm{ed}})$.
Consequently the collar energy of the interpolated logarithms is
at most
\[
 C\bigl(\epsilon_e^2(d_e+F_{\mathrm{ed}})^2+a_{ie}^2\bigr).
\]
Here the cutoff derivative contributes $\rho^{-1}$ and the
Euclidean collar area contributes $O(\rho^2)$. Every term is
bounded by $C\Phi$ after summing incidences: use
\eqref{eq:force-energy-relations}, $\rho=\epsilon_*$ and bounded
$\epsilon_*\ell$. This proves the instantaneous energy bound.
Integrating \eqref{eq:finite-energy} proves the second bound of
\eqref{eq:reconstruction-energy}.

For \eqref{eq:reconstruction-lipschitz}, the endpoint/logarithm
maps have bounded derivatives on the finite ball. Component Green
functions have bounded $L^2$ norms and coefficients $\epsilon_e$.
Their supremum on the truncated component is bounded after
multiplication by $\epsilon_e$. On middle nodes
$0\le\theta_e\le1$, so geodesic endpoint dependence is bounded;
their total area is $O(\rho^2)$. The same smooth endpoint
dependence and the bounded collar logarithms apply to
\eqref{eq:collar-metric}. Integrating along a finite geodesic from
$H$ to $H'$ and using Cauchy--Schwarz proves the stated estimate.
\end{proof}
